\documentclass[11pt,a4paper]{article}

% Packages
\usepackage[utf8]{inputenc}
\usepackage[T1]{fontenc}
\usepackage{amsmath,amssymb,amsthm}
\usepackage{graphicx}
\usepackage{booktabs}
\usepackage{array}
\usepackage{longtable}
\usepackage{hyperref}
\usepackage[margin=1in]{geometry}
\usepackage{authblk}
\usepackage{abstract}
\usepackage{titling}
\usepackage{enumitem}
\usepackage{xcolor}

% Hyperref setup
\hypersetup{
    colorlinks=true,
    linkcolor=blue!50!black,
    citecolor=blue!50!black,
    urlcolor=blue!50!black,
}

% Title and author
\title{From Bell-Pair Dark Matter to Entanglement Cosmogenesis:\\
{\large A Programmatic Conjecture with Testable Predictions}}

\author[1]{ekicyou\thanks{\href{mailto:dot.station@gmail.com}{dot.station@gmail.com}}}
\affil[1]{Independent Researcher, Japan}

\date{May 5, 2026 (Version 1.0)}

\begin{document}

\maketitle

\begin{abstract}
\noindent
We propose a programmatic conjecture for cosmology and particle physics based on a strict \emph{qubit-only ontology}: nothing exists except qubits and their entanglement patterns; concepts such as ``vacuum,'' ``field,'' ``particle,'' and ``spacetime'' are emergent descriptions. Within this framework, dark matter is identified with closed Bell pairs adhering to holographic boundaries, dark energy with the ground-state structure of the boundary register, leptons with open 2-qubit pairs, quarks with 3-qubit triangles, and baryons with 9-qubit \emph{trinity} structures isomorphic to the Shor quantum error-correcting code. We argue that color confinement follows from the no-cloning theorem, that the Pauli exclusion principle emerges from entanglement monogamy combined with the non-existence of AME$(4,2)$ states, and that Hawking radiation is naturally described as discrete qubit ejection events from horizon registers. The framework provides structural answers to several open problems---the cosmological constant problem, the black-hole information paradox, and the lepton/quark asymmetry---and yields approximately thirteen testable quantitative predictions across cosmology, particle physics, and gravitational-wave astronomy. We emphasize that this work is a \emph{programmatic conjecture}, not a completed theory; quantitative derivations and rigorous mathematical formulations are left as open problems for the community. The paper is presented as a dated archive of conceptual structure for potential confrontation with observations on a 5--15 year timescale.
\end{abstract}

\section{Introduction}

\subsection{Motivation: The Crisis of $\Lambda$CDM}

As of 2026, the standard cosmological model ($\Lambda$CDM) faces several persistent observational tensions:
\begin{itemize}[leftmargin=*]
\item The Hubble tension, with disagreement between local and CMB-inferred values of $H_0$ exceeding $6\sigma$~\cite{Riess2022,Planck2020}.
\item DESI 2024--2025 results suggesting dynamical dark-energy evolution rather than a pure cosmological constant~\cite{DESI2024}.
\item JWST observations of unexpectedly mature galaxies at high redshift~\cite{Boylan2023}.
\item More than two decades of null results in direct dark-matter detection experiments~\cite{LUX2017,XENON2023}.
\end{itemize}

These tensions motivate alternative frameworks. Here we explore one: a strict information-theoretic ontology in which all of physics emerges from qubit entanglement.

\subsection{Approach}

We adopt a single ontological principle: \textbf{only qubits and their entanglement patterns exist}. All other concepts---fields, particles, spacetime, vacuum---are emergent descriptions. From this principle we derive structural constraints (qubit number conservation, monogamy-based exclusion, AME-based dimensional limits) and from these constraints we attempt to recover the broad structure of standard physics, with several testable departures.

This approach is in the spirit of Wheeler's ``It from Bit''~\cite{Wheeler1990} but pushed further: not only is information physical, but information \emph{is the only thing that is physical}. The framework also extends ideas from Verlinde's emergent gravity~\cite{Verlinde2010,Verlinde2016}, Wen's string-net liquid~\cite{Wen2017}, and Szangolies's qubit-based gauge group emergence~\cite{Szangolies2025}.

\subsection{Scope and Limitations}

This work is presented as a \emph{programmatic conjecture}, not a completed theory:
\begin{itemize}[leftmargin=*]
\item No claim of mathematical completeness is made.
\item Many quantitative results remain at the level of order-of-magnitude estimates or proposed derivation paths.
\item The author is an independent researcher without academic affiliation; this document is offered for community evaluation rather than to establish authorial priority.
\item The aim is to seed conceptual structure for further development; refinement, refutation, or extension by the broader research community is welcomed.
\end{itemize}

\subsection{Organization}

Section~\ref{sec:foundations} establishes the qubit-only ontology and three strict conservation laws. Section~\ref{sec:hierarchy} derives a particle hierarchy from qubit structures. Section~\ref{sec:dimensions} discusses dimensional emergence and the AME$(4,2)$ obstruction. Sections~\ref{sec:dynamics}--\ref{sec:bh} treat dynamics, cosmology, and black holes. Section~\ref{sec:predictions} catalogs testable predictions. Sections~\ref{sec:related}--\ref{sec:conclusion} discuss related work, open problems, and conclusions.

\section{Foundations: Qubit-Only Ontology}
\label{sec:foundations}

\subsection{The Single Principle}

Our entire framework rests on one ontological commitment:
\begin{quote}
\textit{There exist only qubits and their entanglement patterns. All other concepts are emergent.}
\end{quote}

This minimalist stance excludes: vacuum as an independent substance, fields as fundamental entities, ``creation from nothing,'' and free creation/annihilation of qubits. It commits us to: a fixed-dimensional Hilbert space (modulo cosmological horizon expansion), strictly unitary evolution, and emergent treatment of all other physical concepts.

\subsection{Three Strict Conservation Laws}

From the qubit-only ontology, three strict conservation laws follow:

\paragraph{Conservation A (Qubit Number).}
\begin{equation}
N_{\text{total}} = \frac{A_{\text{cosmic}}}{4\ell_P^2 \ln 2} = \text{constant},
\end{equation}
where $A_{\text{cosmic}}$ is the area of the cosmological horizon. $N_{\text{total}}$ increases discretely only as the cosmological horizon expands.

\paragraph{Conservation B (Total Energy).}
\begin{equation}
E_{\text{total}} = \sum_{i=1}^{N_{\text{total}}} E_i = \text{constant},
\end{equation}
summed over all qubits. Local processes redistribute energy; the total is invariant.

\paragraph{Conservation C (Entanglement Structure).}
The total of all monogamy constraints (in the sense of Coffman-Kundu-Wootters~\cite{CKW2000}) is invariant under local rearrangement.

These conservation laws follow from standard quantum mechanical unitarity but are made explicit as foundational constraints of the framework.

\subsection{Reinterpretation of ``Vacuum''}

In standard quantum field theory, ``vacuum'' is treated as a distinct ontological entity (the ground state of fields). Within our framework:
\begin{quote}
\textit{``Vacuum'' is the collective designation for the state in which all qubits occupy their ground state. It is not a separate entity.}
\end{quote}

Concrete consequences:
\begin{itemize}[leftmargin=*]
\item Vacuum energy = sum of qubit ground-state energies.
\item ``Virtual particles'' = transient excitations of specific qubits (the qubits themselves are real; their excited states are short-lived).
\item Vacuum polarization = local distortion of qubit configurations near charged qubits.
\item Casimir effect = restriction of qubit excitation modes by boundary conditions.
\item Higgs vacuum expectation value = collective ground-state synchronization pattern.
\end{itemize}

The phrase ``borrowing energy from the vacuum'' is replaced by: \emph{exciting a specific qubit}; the lender and borrower are the same qubit.

\subsection{Boundary Registers}

By the holographic principle~\cite{tHooft1993,Susskind1995}, every qubit belongs to some \emph{boundary register} associated with a horizon. Horizons are hierarchical:
\begin{itemize}[leftmargin=*]
\item Each black hole's event horizon
\item The cosmological (de Sitter) horizon of the observable universe
\item Local causal horizons (Rindler-type)
\end{itemize}

The total qubit count of the universe is the sum of qubits across all such boundary registers.

A boundary register has two distinct aspects:

\begin{table}[h]
\centering
\begin{tabular}{lll}
\toprule
\textbf{Aspect} & \textbf{Property} & \textbf{Fluctuation} \\
\midrule
Capacity (qubit count) & Geometric invariant & Deterministic, no fluctuation \\
Content (qubit states) & Quantum, dynamical & Quantum fluctuations \\
\bottomrule
\end{tabular}
\caption{Two aspects of boundary registers.}
\end{table}

\subsection{Boundary Register Temperature}

Each horizon possesses a characteristic temperature:
\begin{align}
T_H &= \frac{\hbar c^3}{8\pi G M k_B} \quad \text{(Schwarzschild)} \\
T_{dS} &= \frac{\hbar H}{2\pi k_B c} \quad \text{(de Sitter)}
\end{align}

We make a key identification:
\begin{equation}
\boxed{T_{\text{register}} = T_{\text{Hawking}}}
\end{equation}

The boundary register is in a thermal state at the Hawking temperature of its horizon. From this single identification follow:
\begin{itemize}[leftmargin=*]
\item Hawking radiation rate (Boltzmann factor)
\item Bekenstein-Hawking entropy formula
\item Black-body spectrum of Hawking radiation
\item Thermal flow between hierarchical horizons
\end{itemize}

\section{Particle Hierarchy from Qubit Structures}
\label{sec:hierarchy}

\subsection{Classification by Qubit Count}

We propose the following identification of particles with qubit structures:

\begin{table}[h]
\centering
\small
\begin{tabular}{lllll}
\toprule
\textbf{Structure} & \textbf{State} & \textbf{Particle} & \textbf{Color} & \textbf{Size} \\
\midrule
1 qubit & free & gauge boson ($\gamma$, $g$) & gluon only & undefined \\
2 qubits & closed (Bell) & Dark Matter & none & boundary area \\
2 qubits & open type A & charged lepton ($e$, $\mu$, $\tau$) & structurally impossible & undefined (point-like) \\
2 qubits & open type B & neutrino ($\nu$) & structurally impossible & undefined (point-like) \\
3 qubits & open (triangle) & quark & exposed & cannot exist alone \\
3 qubits & bulk-projected & W, Z bosons & none & short-range \\
6 qubits & $3 + \bar{3}$ & meson & singlet & hadron-scale \\
\textbf{9 qubits} & \textbf{trinity (Shor-isomorphic)} & \textbf{baryon} & \textbf{singlet} & \textbf{proton-scale} \\
many qubits & hierarchical & nuclei, atoms, molecules & singlet & various \\
\bottomrule
\end{tabular}
\caption{Particle hierarchy from qubit structures.}
\label{tab:hierarchy}
\end{table}

\subsection{Lepton/Quark Asymmetry as Structural Necessity}

The Standard Model's long-standing puzzle---\textit{why do quarks alone carry color?}---admits a structural answer:
\begin{quote}
\textit{Leptons are 2-qubit structures; quarks are 3-qubit structures. The difference $2 \neq 3$ directly produces the presence or absence of color.}
\end{quote}

Specifically, color charge is carried by 3-qubit triangle structures. Since 2-qubit structures contain no triangle, leptons are \emph{structurally incapable} of carrying color charge or participating in strong interactions. This elevates SU(3) non-coupling for leptons from a phenomenological fact to a structural necessity.

\subsection{The 9-Qubit Trinity = Baryon}

The mechanism by which 3-qubit quarks combine into 9-qubit baryons follows from quantum-information principles. Five independent considerations converge on the same conclusion:

\begin{table}[h]
\centering
\small
\begin{tabular}{lll}
\toprule
\textbf{Perspective} & \textbf{Single 3-qubit (unstable)} & \textbf{9-qubit trinity (stable)} \\
\midrule
Topology & open disk ($\chi = 1$) & closed sphere ($\chi = 2$) \\
Entanglement saturation & 3 exposed edges & all 21 edges shared \\
Statistics & forced anyon collapse & complex consistency \\
Quantum information & Shor-block extraction forbidden & complete Shor code \\
AME hierarchy & blocked at AME(4,2) & pure volumetric realization \\
\bottomrule
\end{tabular}
\caption{Five-fold convergence on 9-qubit trinity stability.}
\end{table}

\paragraph{Shor-code isomorphism.}
The Shor 9-qubit quantum error-correcting code~\cite{Shor1995} has the structure
\begin{align}
|\bar{0}\rangle &= \frac{1}{2\sqrt{2}}\bigotimes_{i=1}^{3}(|000\rangle + |111\rangle)_i \\
|\bar{1}\rangle &= \frac{1}{2\sqrt{2}}\bigotimes_{i=1}^{3}(|000\rangle - |111\rangle)_i
\end{align}
which is the hierarchical composition of three GHZ-like states (each on 3 qubits).

\paragraph{Color confinement from no-cloning.}
A crucial consequence: extracting a single 3-qubit block from a Shor-encoded state corresponds to extracting partial information from a quantum error-correcting code, which is forbidden by the no-cloning theorem~\cite{Wootters1982}. Therefore:
\begin{quote}
\textit{Color confinement is structurally derivable from the no-cloning theorem.}
\end{quote}

The geometric realization is the \emph{triaugmented triangular prism}: 9 vertices, 21 edges, 14 triangular faces, Euler characteristic $\chi = 2$, with every edge shared between exactly two faces (complete saturation).

\subsection{Open 2-Qubit Pair = Lepton}

An open 2-qubit pair (with partial internal entanglement and external entanglement capacity) is identified with a lepton.

\paragraph{Size-undefined nature of the electron.}
Three considerations give the electron's apparent point-like nature ($r_e < 10^{-18}$ m experimentally) a structural origin:
\begin{enumerate}[leftmargin=*]
\item \textbf{No closed surface}: a 2-qubit structure has no closed bounding surface in 3D.
\item \textbf{Not boundary-adhered}: unlike Bell pairs, open 2-qubit pairs float in bulk.
\item \textbf{Dynamical entanglement}: continuous environmental re-entanglement precludes a fixed shape.
\end{enumerate}

The electron is thus not ``a point-like particle'' but ``a structure to which the concept of size does not apply.''

\paragraph{Five observed properties.}
\begin{table}[h]
\centering
\begin{tabular}{ll}
\toprule
\textbf{Observation} & \textbf{Structural origin} \\
\midrule
Spin-$\tfrac{1}{2}$ & open 2-qubit minimal half-integer spin carrier \\
Electric charge & residual U(1) coupling from monogamy excess \\
Mass $m_e \ll m_{\rm DM}$ & partial internal vibration $<$ complete vibration \\
No color & absence of triangle structure \\
Weak coupling & dynamical coupling to 3-qubit (W, Z) structures \\
\bottomrule
\end{tabular}
\end{table}

\section{Dimensional Emergence}
\label{sec:dimensions}

\subsection{The Qubit Ladder}

The dimensionality of spacetime emerges in stages corresponding to qubit count:

\begin{table}[h]
\centering
\small
\begin{tabular}{lllll}
\toprule
\textbf{Qubits} & \textbf{Simplex} & \textbf{Shape} & \textbf{Dimension emerged} & \textbf{BH stage} \\
\midrule
1 & 0-simplex & point & 0D & --- \\
2 & 1-simplex & edge & 1D = pre-spatial & proto-BH (DM) \\
3 & 2-simplex & triangle & 2D = boundary surface & semi-complete BH \\
4 & 3-simplex & tetrahedron & 3D, but blocked by AME$(4,2)$ & --- \\
\textbf{9} & \textbf{trinity} & \textbf{triaug. prism} & \textbf{3D (volume)} & \textbf{minimal complete BH} \\
\bottomrule
\end{tabular}
\caption{Dimensional emergence through qubit ladder.}
\end{table}

\subsection{The AME$(4,2)$ Obstruction}

The Higuchi-Sudbery theorem~\cite{Higuchi2000} establishes that no Absolutely Maximally Entangled state on 4 qubits in dimension 2 exists. This implies:
\begin{itemize}[leftmargin=*]
\item ``Perfect 4-qubit tetrahedra'' are mathematically forbidden.
\item 3D volume cannot be perfectly realized at 4-qubit scale.
\item This obstruction provides an intrinsic origin for quantum-gravitational noise in 3D space.
\end{itemize}

The AME$(n,2)$ existence pattern for small $n$:
\begin{center}
\begin{tabular}{ccccccc}
\toprule
$n$ & 2 & 3 & \textbf{4} & 5 & 6 & 7 \\
\midrule
AME exists & \checkmark & \checkmark & $\boldsymbol{\times}$ & \checkmark & \checkmark & $\times$ \\
\bottomrule
\end{tabular}
\end{center}

The framework requires skipping over $n=4$ to $n=5$ (open 2+3 composite, intermediate states) or $n=9$ (3:3:3 trinity, primary route).

\subsection{Two-Stage Activation of the Pauli Exclusion Principle}

Corresponding to dimensional emergence, the Pauli exclusion principle activates in two stages:

\begin{description}
\item[Phase 3-2 (3-qubit triangle, 2D): proto-exclusion] Antisymmetric components in Young-tableau decomposition; anyonic statistics permitted.
\item[Phase 3-3 (9-qubit trinity, 3D volume): standard Pauli exclusion] Spin-statistics theorem holds; fermion/boson dichotomy fixed.
\end{description}

The roots of the exclusion principle reside in the qubit layer itself:
\begin{enumerate}[leftmargin=*]
\item Entanglement monogamy (CKW inequality)
\item AME$(4,2)$ non-existence
\item Singlet antisymmetry under exchange
\end{enumerate}

These pre-existing repulsion structures are projected onto spacetime as the observed Pauli exclusion principle.

\section{Dynamics}
\label{sec:dynamics}

\subsection{Reaction Cascade}

The central dynamics consist of three hierarchical reactions:
\begin{align}
\text{Stage 1:} \quad & 2\text{-qubit} + 1\text{-qubit} \rightleftharpoons 3\text{-qubit} \quad (T_c \sim 10^{16}~\text{GeV}) \\
\text{Stage 2:} \quad & 3 \times 3\text{-qubit} \rightleftharpoons 9\text{-qubit trinity} \quad (T_{\rm QCD} \sim 200~\text{MeV}) \\
\text{Stage 3:} \quad & 3\text{-qubit} \rightleftharpoons 2\text{-qubit} + 1\text{-qubit} \quad \text{(strongly suppressed)}
\end{align}

Temperature freeze-out of these reactions determines cosmological density ratios $\Omega_{\rm DM}$, $\Omega_b$, $\Omega_\gamma$.

\subsection{Topological Boundary Transitions}

Open$\leftrightarrow$closed transitions of qubit pairs are subject to a special constraint:
\begin{quote}
\textit{Open$\leftrightarrow$closed qubit-pair transitions are forbidden in bulk; they must proceed via the boundary.}
\end{quote}

This implies that closed Bell pairs (DM) cannot appear directly as products of bulk reactions; their formation requires boundary mediation.

\subsection{Boundary-Bulk Qubit Exchange}

Major matter-related processes are described as qubit exchanges between boundary registers and bulk:
\begin{table}[h]
\centering
\begin{tabular}{ll}
\toprule
\textbf{Process} & \textbf{Qubit pathway} \\
\midrule
Matter formation & boundary register $\to$ bulk projection \\
Pair annihilation & bulk particle qubits $\to$ boundary register + photons \\
Pair creation & boundary register activation + photon absorption $\to$ bulk \\
BH formation & bulk matter qubits $\to$ horizon register \\
Hawking radiation & horizon register $\to$ bulk emission \\
\bottomrule
\end{tabular}
\end{table}

In all processes, the total qubit number is strictly conserved.

\section{Cosmology}
\label{sec:cosmology}

\subsection{Cosmic History Scenario}

\begin{enumerate}[leftmargin=*]
\item \textbf{Phase 1 (pre-spatial):} Only qubit structures exist; no space, time, or symmetry.
\item \textbf{Phase 2 (dynamic equilibrium):} Stage-1 reaction in dynamic equilibrium at Planck energies.
\item \textbf{Phase 3-1 (2D emergence)} at $T_c \sim 10^{16}$ GeV: 3-qubit dominance; 2D boundary surface emerges; proto-exclusion activates.
\item \textbf{Phase 3-2 (``2D era''):} Tessellation by 3-qubit triangles completes the holographic boundary. This phase is consistent with CDT dimensional reduction observations~\cite{Ambjorn2005}.
\item \textbf{Phase 3-3 (3D emergence):} 9-qubit trinity formation freezes Stage 2; true 3D volume; standard Pauli exclusion fixed.
\item \textbf{Phase 4 (vibrational inflation):} Distance-dependent exclusion pressure drives oscillatory inflation; cumulative expansion $\sim$60 e-folds; oscillation damping = reheating.
\item \textbf{Phase 5 (vacuum structure established):} Boundary register settles; multi-qubit structures (baryons, nuclei) form from successful Stage-1 transitions; failed transitions remain as DM-precursor Bell pairs.
\item \textbf{Phase 6 (matter-dominated era):} Hot DM clusters form; Bell-pair register fluctuations observed as cold DM.
\item \textbf{Phase 7 (dark-energy era / present):} $\Omega_{\rm DE} \approx 0.68$ from Bell-pair register average density; $\Omega_{\rm DM} \approx 0.27$ from fluctuations and excited clusters.
\item \textbf{Phase 8 (BH formation / matter consumption):} Falling matter and pair annihilation introduce qubits to boundary registers.
\item \textbf{Phase 9 (de Sitter asymptote):} All matter incorporated into boundary registers; only registers remain.
\end{enumerate}

\subsection{Vibrational Inflation}

Pressure from distance-dependent exclusion (Fermi-gas-like):
\begin{equation}
P = \frac{2}{5} n E_F, \quad E_F \propto n^{2/3}, \quad P_{\rm internal} \propto n^{5/3}
\end{equation}

The effective pressure $P_{\rm eff} = P_{\rm internal} - P_{\rm external}$ changes sign with density, driving oscillation:
\begin{itemize}[leftmargin=*]
\item Phase 4a: high density, strong positive pressure $\to$ rapid expansion
\item Phase 4b: inertial overshoot of equilibrium distance
\item Phase 4c: negative pressure $\to$ recoil contraction
\item Phase 4d: damped oscillation; energy dissipation per cycle
\item Phase 4e: oscillation termination = reheating
\end{itemize}

This picture has thematic affinity with bouncing/cyclic cosmology~\cite{Steinhardt2002,LQC} but proceeds via discrete oscillation rather than continuous bounce.

\subsection{Unified Description of the Dark Sector}

The boundary register in different states gives rise to all ``dark'' phenomena:
\begin{table}[h]
\centering
\begin{tabular}{lll}
\toprule
\textbf{Register state} & \textbf{Observed phenomenon} & \textbf{Contribution} \\
\midrule
Uniform ground state & Dark Energy (DE) & $\Omega_{\rm DE} \approx 0.68$ \\
Ground-state local fluctuations & cold DM (axion-like) & part of $\Omega_{\rm DM}$ \\
Excited clusters & hot DM & galactic halos \\
Collective vibrational modes & Higgs boson & mass attribution \\
Sync with register & particle masses & $mc^2 = \hbar\omega_{\rm sync}$ \\
\bottomrule
\end{tabular}
\end{table}

\subsection{Resolution of the Cosmological Constant Problem}

The standard QFT prediction
\begin{equation}
\rho_{\rm vac}^{\rm QFT} \sim M_{\rm Planck}^4 \sim 10^{120} \rho_{\rm DE}^{\rm obs}
\end{equation}
assumes \emph{independent} contributions from each fundamental field. In our framework:
\begin{quote}
\textit{All ``fields'' are emergent excitation modes of the single qubit register. Zero-point energies cannot be summed independently; the total is given by a single sum over the register, naturally yielding a small value.}
\end{quote}

A complementary thermal explanation: the cosmological horizon temperature $T_{dS} \sim 10^{-30}$ K places the boundary register essentially in its ground state, with Boltzmann-suppressed excitations. Therefore $\rho_{\rm DE}$ is naturally small.

\section{Black Holes and Information}
\label{sec:bh}

\subsection{Black Hole as Condensed Boundary Register}

A black hole is interpreted as a region where bulk matter has been incorporated into a boundary register:
\begin{itemize}[leftmargin=*]
\item BH formation = quantization of falling matter into horizon register at the event horizon.
\item BH ``interior'' is materially absent; the horizon is everything.
\item Singularities do not form (no matter reaches them).
\item Quantum gravity is unnecessary (spacetime terminates at the horizon level).
\end{itemize}

This aligns with AMPS firewall considerations~\cite{AMPS2013}, Mathur's fuzzball proposal~\cite{Mathur2005}, and the broader emergent-gravity program.

\subsection{Hawking Radiation as Qubit Ejection}

In our framework, Hawking radiation is described as discrete events at the horizon:
\begin{enumerate}[leftmargin=*]
\item Two adjacent horizon qubits, both in ground state.
\item Local qubit-qubit interaction excites both.
\item Asymmetric energy distribution: one qubit returns to ground state (releasing energy); the other absorbs and exceeds horizon binding.
\item Result: one qubit remains on horizon, one ejects to bulk as a single-qubit photon. Horizon area decreases by one Planck plaquette.
\end{enumerate}

This is equivalent to the standard QFT virtual-pair picture but recasts ``negative-energy partner'' as ``qubit returning to ground state,'' avoiding philosophical issues with negative-energy particles.

\subsection{Automatic Resolution of the Information Paradox}

Each ejected qubit was previously entangled with remaining horizon qubits. As evaporation progresses, entanglement is transferred outward. The Page curve~\cite{Page1993} arises automatically: radiation entanglement entropy peaks at half-evaporation. By complete evaporation, all original information has been transferred to the external Hawking radiation.

\begin{quote}
\textit{Hawking radiation = stepwise readout of black-hole information.}
\end{quote}

No firewall is required; no information is lost; unitarity is preserved by construction.

\subsection{Unification of Four Major Processes}

The framework unifies four processes under a single mechanism (boundary-bulk qubit exchange):

\begin{table}[h]
\centering
\begin{tabular}{lll}
\toprule
\textbf{Process} & \textbf{Direction} & \textbf{Pathway} \\
\midrule
BH formation & bulk $\to$ horizon & matter qubits absorbed by register \\
Hawking radiation & horizon $\to$ bulk & register qubits ejected \\
Pair annihilation & bulk $\to$ boundary & particle qubits absorbed + photons emitted \\
Pair creation & boundary $\to$ bulk & register activated + photons absorbed \\
\bottomrule
\end{tabular}
\end{table}

All processes preserve total qubit number, total energy, and unitarity by construction.

\subsection{Hawking Thermodynamics from Register Statistics}

The identification $T_{\rm register} = T_{\rm Hawking}$ yields:

\paragraph{Hawking radiation rate (Boltzmann-suppressed):}
\begin{equation}
\text{rate} \propto \exp(-E_{\rm eject}/k_B T_H)
\end{equation}

\paragraph{Black-body spectrum:} Register at $T_H$ in thermal equilibrium $\to$ thermal spectrum of ejected qubits.

\paragraph{Bekenstein-Hawking entropy:}
\begin{equation}
S_{\rm BH} = \frac{A}{4\ell_P^2}
\end{equation}
arises as thermal entropy of the register at $T_H$.

\paragraph{Inter-horizon thermal flow:} Temperature difference between BH horizon ($T_H \sim 10^{-7}$ K for stellar mass) and cosmological horizon ($T_{dS} \sim 10^{-30}$ K) drives Hawking radiation as heat flow from high-temperature to low-temperature register.

\section{Testable Predictions}
\label{sec:predictions}

\subsection{Quantitative Predictions}

\begin{longtable}{p{4.5cm}p{3.5cm}p{4cm}p{2cm}}
\toprule
\textbf{Observable} & \textbf{Predicted value} & \textbf{Verification} & \textbf{Timeline} \\
\midrule
Inflation temperature $T_c$ & $\sim 10^{16}$ GeV & CMB B-mode, primordial GW & 2032+ \\
Tensor-to-scalar ratio $r$ & 0.001--0.01 & LiteBIRD & 2032+ \\
CMB $\mu$-distortion & $\sim 10^{-7}$ & PIXIE/PRISM-class missions & proposed \\
Primordial GW peak frequency & $10^{-9}$--$10^{-7}$ Hz & PTA (NANOGrav, etc.) & ongoing \\
CMB power-spectrum oscillations & specific peaks/dips & Planck + successors & verifiable \\
Primordial GW oscillation traces & discrete frequency structure & NANOGrav, LISA & ongoing \\
$\Omega_{\rm DM} \approx 0.27$ & cascade freeze-out & Planck, DESI consistency & existing \\
$\Omega_{\rm DE} \approx 0.68$ & register ground state density & Planck, DESI & existing \\
DM/baryon $\approx 5.4$ & freeze-out residual ratio & same as above & existing \\
DM mass distribution & bimodal (hot + cold DM) & direct-detection mass scans & long-term \\
511 keV photon total energy & $1022$ keV with $< 1$ eV deviation & topological transition test & existing apparatus \\
Galactic-center DM density & micro-enhancement proportional to annihilation rate & gamma-ray excess observations & ongoing \\
GZK-cutoff anomaly & partial explanation via qubit ejection & cosmic-ray observatories & ongoing \\
\bottomrule
\caption{Quantitative testable predictions.}
\end{longtable}

\subsection{Qualitative Predictions}

\begin{enumerate}[leftmargin=*]
\item Direct DM detection will remain perpetually null (boundary adherence).
\item DM halo shape correlates with rotation (boundary Kerr-like rotation).
\item Large-scale structure exhibits subtle anisotropy (3-qubit transition initial conditions).
\item Quantum gravity research is misdirected (no spacetime exists at Planck scale).
\item No matter reaches BH interior (qubit-ization at horizon, no singularity).
\item Gravitational-wave B-mode shows characteristic spectral tilt (oscillatory transition).
\item All interactions describable in entanglement-mode language (consistent with Wen's string-net framework).
\item Leptons cannot acquire color charge (structural).
\item Quarks cannot be isolated (Shor-block extraction = no-cloning violation).
\item Pair annihilation isomorphic to micro-scale BH formation/evaporation.
\item Bullet Cluster constraint automatically satisfied (2+2 volume formation forbidden $\to$ DM self-interaction zero).
\item Qubit-ejection effects at high energy are Boltzmann-suppressed at all currently accessible energies.
\item Hawking radiation automatically preserves information (Page curve, no firewall).
\end{enumerate}

\subsection{Verification Strategy: Gravitational Probes of Pre-Inflationary Physics}

\begin{table}[h]
\centering
\begin{tabular}{lll}
\toprule
\textbf{Channel} & \textbf{Probe} & \textbf{Era explored} \\
\midrule
Primordial GW (B-mode) & LiteBIRD (2032), CMB-S4 & end of inflation \\
Stochastic GW background & PTA, LISA (2035), DECIGO & post-inflation reheating \\
CMB spectral distortion & PIXIE/PRISM & $z \approx 10^6$ \\
Cosmic neutrino background & PTOLEMY etc. & $z \approx 10^9$ \\
\bottomrule
\end{tabular}
\end{table}

CMB $\mu$-distortion at $z \approx 10^6$ is particularly powerful as a direct probe of the end of inflation.

\section{Internal Consistency}

\subsection{Five-Path Convergence on the Quantum Matter Scale}

The proton size ($\sim 10^{-15}$ m) emerges from five independent paths within the framework:
\begin{itemize}[leftmargin=*]
\item \textbf{Path A}: holographic degree-of-freedom volumization
\item \textbf{Path B}: exclusion-principle action range
\item \textbf{Path C}: exclusion equilibrium distance = volume/area ratio
\item \textbf{Path D}: triaugmented prism size of 9-qubit trinity
\item \textbf{Path E}: Bell-pair register local fluctuation wavelength
\end{itemize}

The convergence is taken as evidence of internal consistency.

\subsection{Five-Perspective Convergence on Trinity Stability}

As detailed in Section~\ref{sec:hierarchy}, five perspectives (topology, entanglement saturation, statistics, quantum information, AME hierarchy) all converge on the conclusion that single 3-qubit structures are unstable while 9-qubit trinities are stable.

\section{Comparison with Existing Frameworks}
\label{sec:related}

This work draws on and extends several existing programs:

\paragraph{Field-theoretic DM candidates.} Sterino model~\cite{Krolikowski2007}, Cooper-pair DM~\cite{Alexander2024}, neutrino superfluidity~\cite{Kapusta2004}, singlet scalar DM~\cite{Burgess2001}.

\paragraph{Holographic DM origins.} Holographic DM~\cite{Fichet2026}, bulk-BH cosmological DM~\cite{Bulk2022}, emergent DM on holographic screens~\cite{Emergent2017}, DM from holography via CKN bound~\cite{Holo2025}.

\paragraph{Quantum-information spacetime emergence.} Szangolies~\cite{Szangolies2025}, Wen's string-net liquid~\cite{Wen2017}, Van Raamsdonk's entanglement-spacetime correspondence~\cite{VanRaamsdonk2010}, ER=EPR~\cite{Maldacena2013}, CKN bound~\cite{Cohen1999}.

\paragraph{Emergent gravity / quantum gravity rejection.} Verlinde~\cite{Verlinde2010,Verlinde2016}, Jacobson~\cite{Jacobson1995}, Padmanabhan~\cite{Padmanabhan2010}, Sakharov~\cite{Sakharov1967}, AMPS firewall paradox~\cite{AMPS2013}, Mathur fuzzball~\cite{Mathur2005}.

\paragraph{Bouncing/cyclic cosmology.} Ekpyrotic/cyclic universe~\cite{Steinhardt2002}, loop quantum cosmology~\cite{LQC}, matter-bounce inflation, quintessential inflation.

\paragraph{Quantum error correction and physics.} Shor code~\cite{Shor1995}, holographic codes (HaPPY)~\cite{Pastawski2015}.

The novelty of the present work is the \emph{integrated synthesis}: not the individual elements (most of which exist in the literature), but the unified framework that organizes them under a single qubit-only ontology.

\section{Open Problems and Future Work}
\label{sec:open}

\subsection{Quantitative Challenges}

The following derivations remain incomplete:
\begin{itemize}[leftmargin=*]
\item Cascade equilibrium constants and the physical determination of each $\Delta E$.
\item Thermal-relic-style calculation of $\Omega_{\rm DM}$, $\Omega_{\rm DE}$, $\Omega_b$.
\item DM/baryon ratio from freeze-out temperatures.
\item Determination of vibration amplitude and period in inflation.
\item Three-generation lepton mass ratios from open-2-qubit mode hierarchy.
\item Quantitative qubit-ejection rate and its energy dependence.
\item Higgs vev (246 GeV) from register vibrational coupling scale.
\end{itemize}

\subsection{Internal Consistency Issues}

\begin{itemize}[leftmargin=*]
\item Mathematical formalism of the boundary register (Hilbert space construction, Hamiltonian).
\item Uniqueness of the triaugmented prism realization for the 9-qubit trinity.
\item Mathematical formulation of three-generation discrete mode hierarchy.
\item Mechanism for excited-Bell-pair generation and decay.
\end{itemize}

\subsection{Observational Consistency}

\begin{itemize}[leftmargin=*]
\item $r < 0.036$ constraint excludes Planck-scale direct inflation (consistent with $T_c \sim 10^{16}$ GeV).
\item $N_{\rm eff}$ and structure-formation impacts require detailed computation.
\item CMB power-spectrum oscillation predictions need preliminary verification against existing Planck data.
\item 511 keV photon energy at $<$ eV precision: feasibility of direct test of topological transition rule.
\end{itemize}

\subsection{Literature}

A thorough review of recent works (Szangolies 2025, Fichet 2026, Cooper-pair DM 2024, holographic codes) is needed to refine the boundary between this work and existing literature.

\section{Discussion}

\subsection{Philosophical Position}

If the framework is correct, the following holds:
\begin{quote}
\textit{Physics is the organization of information.}
\end{quote}

This is a radicalization of Wheeler's ``It from Bit'': not only do informational principles underlie physics, but \emph{nothing exists except information patterns}. There is some affinity to Leibniz's monadology, but recast in modern quantum-information terms.

\subsection{Position Relative to Mainstream Physics}

This framework does not contradict standard physics; rather, it proposes that standard QFT is the \emph{effective field theory} of the underlying qubit register. Specifically:
\begin{itemize}[leftmargin=*]
\item All standard-model predictions follow automatically as low-energy consequences.
\item Deviations would appear only at extreme energies (Planck-scale) or in cosmological observations.
\item The cosmological constant problem and the BH information paradox are resolved structurally rather than via additional fine-tuning or new physics.
\end{itemize}

The framework is therefore positioned as a \emph{complement} to mainstream physics, not a replacement.

\section{Conclusion}
\label{sec:conclusion}

We have presented a programmatic conjecture for cosmology and particle physics based on a strict qubit-only ontology. The framework:

\begin{enumerate}[leftmargin=*]
\item Identifies dark matter, dark energy, Higgs phenomenology, and particle masses as different states of a unified Bell-pair boundary register.
\item Derives lepton/quark asymmetry, color confinement, and the Pauli exclusion principle from quantum-information principles.
\item Recasts Hawking radiation as discrete qubit-ejection events, automatically resolving the information paradox.
\item Proposes structural resolutions to the cosmological constant problem and several open issues in $\Lambda$CDM.
\item Yields approximately thirteen testable quantitative predictions and many qualitative ones.
\end{enumerate}

The framework is not a completed theory. Quantitative derivations and rigorous mathematical formulations remain open. We submit this document as a dated archive of conceptual structure, for evaluation, refinement, refutation, or extension by the broader research community.

We expect that observations on a 5--15 year timescale---LiteBIRD, LISA, PIXIE-class missions, ongoing PTA campaigns, the Vera C. Rubin Observatory, and high-precision pair-annihilation measurements---will provide opportunities to test specific predictions of this framework.

\section*{Acknowledgments}

The author gratefully acknowledges extensive conceptual development through dialogue with Claude Opus 4.7 (Anthropic, San Francisco, CA, USA). The AI system served as a sparring partner for hypothesis refinement, literature cross-referencing, and structural consistency checks. All scientific claims, errors, and final positions are the author's responsibility.

\section*{Statement of Purpose}

The author is an independent researcher without formal academic affiliation. This document is submitted as a programmatic conjecture rather than a claim of completed theory. The aim is to seed ideas for community development, not to establish authorial priority. The author welcomes refinement, refutation, and extension by the broader research community.

\section*{License}

This document is released under Creative Commons Attribution 4.0 International (CC-BY 4.0). Free use, modification, and distribution are permitted with attribution.

\begin{thebibliography}{99}

\bibitem{Wheeler1990} J.~A.~Wheeler, ``Information, physics, quantum: The search for links,'' in \emph{Complexity, Entropy, and the Physics of Information} (1990).

\bibitem{Bekenstein1973} J.~D.~Bekenstein, ``Black holes and entropy,'' Phys. Rev. D \textbf{7}, 2333 (1973).

\bibitem{Hawking1975} S.~W.~Hawking, ``Particle creation by black holes,'' Commun. Math. Phys. \textbf{43}, 199 (1975).

\bibitem{tHooft1993} G.~'t Hooft, ``Dimensional reduction in quantum gravity,'' arXiv:gr-qc/9310026 (1993).

\bibitem{Susskind1995} L.~Susskind, ``The world as a hologram,'' J. Math. Phys. \textbf{36}, 6377 (1995).

\bibitem{Maldacena1997} J.~Maldacena, ``The large N limit of superconformal field theories and supergravity,'' Adv. Theor. Math. Phys. \textbf{2}, 231 (1998).

\bibitem{Page1993} D.~N.~Page, ``Information in black hole radiation,'' Phys. Rev. Lett. \textbf{71}, 3743 (1993).

\bibitem{Shor1995} P.~W.~Shor, ``Scheme for reducing decoherence in quantum computer memory,'' Phys. Rev. A \textbf{52}, 2493 (1995).

\bibitem{Wootters1982} W.~K.~Wootters and W.~H.~Zurek, ``A single quantum cannot be cloned,'' Nature \textbf{299}, 802 (1982).

\bibitem{CKW2000} V.~Coffman, J.~Kundu, and W.~K.~Wootters, ``Distributed entanglement,'' Phys. Rev. A \textbf{61}, 052306 (2000).

\bibitem{Higuchi2000} A.~Higuchi and A.~Sudbery, ``How entangled can two couples get?'' Phys. Lett. A \textbf{273}, 213 (2000).

\bibitem{Verlinde2010} E.~Verlinde, ``On the origin of gravity and the laws of Newton,'' JHEP \textbf{04}, 029 (2011).

\bibitem{Verlinde2016} E.~Verlinde, ``Emergent gravity and the dark universe,'' SciPost Phys. \textbf{2}, 016 (2017).

\bibitem{Wen2017} X.-G.~Wen, ``A theory of 2+1D bosonic topological orders,'' arXiv:1709.03824.

\bibitem{Szangolies2025} J.~Szangolies, ``Spacetime and standard model from qubits,'' Entropy \textbf{27}, 569 (2025); arXiv:2512.17328.

\bibitem{Jacobson1995} T.~Jacobson, ``Thermodynamics of spacetime: The Einstein equation of state,'' Phys. Rev. Lett. \textbf{75}, 1260 (1995).

\bibitem{Padmanabhan2010} T.~Padmanabhan, ``Thermodynamical aspects of gravity: New insights,'' Rep. Prog. Phys. \textbf{73}, 046901 (2010).

\bibitem{Sakharov1967} A.~D.~Sakharov, ``Vacuum quantum fluctuations in curved space and the theory of gravitation,'' Sov. Phys. Dokl. \textbf{12}, 1040 (1968).

\bibitem{AMPS2013} A.~Almheiri, D.~Marolf, J.~Polchinski, and J.~Sully, ``Black holes: Complementarity or firewalls?'' JHEP \textbf{02}, 062 (2013).

\bibitem{Mathur2005} S.~D.~Mathur, ``The fuzzball proposal for black holes: An elementary review,'' Fortsch. Phys. \textbf{53}, 793 (2005).

\bibitem{Maldacena2013} J.~Maldacena and L.~Susskind, ``Cool horizons for entangled black holes,'' Fortsch. Phys. \textbf{61}, 781 (2013).

\bibitem{VanRaamsdonk2010} M.~Van Raamsdonk, ``Building up spacetime with quantum entanglement,'' Gen. Rel. Grav. \textbf{42}, 2323 (2010).

\bibitem{Cohen1999} A.~G.~Cohen, D.~B.~Kaplan, and A.~E.~Nelson, ``Effective field theory, black holes, and the cosmological constant,'' Phys. Rev. Lett. \textbf{82}, 4971 (1999).

\bibitem{Pastawski2015} F.~Pastawski, B.~Yoshida, D.~Harlow, and J.~Preskill, ``Holographic quantum error-correcting codes,'' JHEP \textbf{06}, 149 (2015).

\bibitem{Fichet2026} S.~Fichet, E.~Megías, and M.~Quirós, ``Holographic dark matter,'' arXiv:2602.13393.

\bibitem{Bulk2022} (anonymous), ``Cosmological dark matter from a bulk black hole,'' arXiv:2212.13268.

\bibitem{Emergent2017} (anonymous), ``Emergent dark matter on the holographic screen,'' arXiv:1712.09326.

\bibitem{Holo2025} (anonymous), ``Dark matter from holography,'' arXiv:2511.10617.

\bibitem{Krolikowski2007} W.~Królikowski, ``Sterino as cold dark matter and its hot lepton counterpart,'' arXiv:0712.0937.

\bibitem{Alexander2024} S.~Alexander, R.~Bernardo, and A.~Gilmer, ``Cooper-pair dark matter,'' arXiv:2405.08874.

\bibitem{Kapusta2004} J.~Kapusta, ``Neutrino superfluidity,'' Phys. Rev. Lett. \textbf{93}, 251801 (2004).

\bibitem{Burgess2001} C.~P.~Burgess, M.~Pospelov, and T.~ter Veldhuis, ``The minimal model of nonbaryonic dark matter: A singlet scalar,'' Nucl. Phys. B \textbf{619}, 709 (2001).

\bibitem{Steinhardt2002} P.~J.~Steinhardt and N.~Turok, ``A cyclic model of the universe,'' Science \textbf{296}, 1436 (2002).

\bibitem{LQC} A.~Ashtekar, T.~Pawlowski, and P.~Singh, ``Quantum nature of the big bang,'' Phys. Rev. Lett. \textbf{96}, 141301 (2006).

\bibitem{Ambjorn2005} J.~Ambj{\o}rn, J.~Jurkiewicz, and R.~Loll, ``Spectral dimension of the universe,'' Phys. Rev. Lett. \textbf{95}, 171301 (2005).

\bibitem{Riess2022} A.~G.~Riess et al., ``A comprehensive measurement of the local value of the Hubble constant,'' Astrophys. J. \textbf{934}, L7 (2022).

\bibitem{Planck2020} Planck Collaboration, ``Planck 2018 results. VI. Cosmological parameters,'' A\&A \textbf{641}, A6 (2020).

\bibitem{DESI2024} DESI Collaboration, ``DESI 2024 VI: Cosmological constraints from the measurements of baryon acoustic oscillations,'' arXiv:2404.03002.

\bibitem{Boylan2023} C.~Boylan-Kolchin, ``Stress testing $\Lambda$CDM with high-redshift galaxy candidates,'' Nat. Astron. \textbf{7}, 731 (2023).

\bibitem{LUX2017} LUX Collaboration, ``Results from a search for dark matter in the complete LUX exposure,'' Phys. Rev. Lett. \textbf{118}, 021303 (2017).

\bibitem{XENON2023} XENON Collaboration, ``First dark matter search with nuclear recoils from the XENONnT experiment,'' Phys. Rev. Lett. \textbf{131}, 041003 (2023).

\end{thebibliography}

\end{document}
